\ifmostrarGabarito
\begin{minted}{python}
class ContaBancaria:
    def __init__(self, titular: str, numero_conta: int, saldo_inicial: float = 0.0):
        self.titular = titular
        self.numero_conta = numero_conta
        self.saldo = float(saldo_inicial)
        
    def depositar(self, valor: float) -> bool:
        if valor > 0:
            self.saldo += valor
            return True
        return False
        
    def sacar(self, valor: float) -> bool:
        if 0 < valor <= self.saldo:
            self.saldo -= valor
            return True
        return False
        
    def __str__(self) -> str:
        return f"Conta {self.numero_conta} - Titular: {self.titular} - Saldo: R$ {self.saldo:.2f}"


class ContaEspecial(ContaBancaria):
    def __init__(self, titular: str, numero_conta: int, saldo_inicial: float = 0.0, limite_cheque_especial: float = 500.0):
        super().__init__(titular, numero_conta, saldo_inicial)
        self.limite = float(limite_cheque_especial)
        
    def sacar(self, valor: float) -> bool:
        if 0 < valor <= (self.saldo + self.limite):
            self.saldo -= valor
            return True
        return False
        
    def __str__(self) -> str:
        base_str = super().__str__()
        return f"{base_str} (Limite: R$ {self.limite:.2f})"
\end{minted}

\textbf{Explicação:}
A classe \texttt{ContaBancaria} estabelece o modelo base com encapsulamento de estado e método especial \texttt{\_\_str\_\_}. A subclasse \texttt{ContaEspecial} reutiliza o construtor da mãe com \texttt{super().\_\_init\_\_}, sobrescreve \texttt{sacar} para incorporar a margem de limite especial e especializa \texttt{\_\_str\_\_} estendendo a representação da superclasse.
\else
\linhasResposta{15}
\clearpage
\linhasResposta[FOLHA DE CONTINUACAO / RASCUNHO -- QUESTAO Q15]{30}
\clearpage
\fi
